%% figures/fig4_prefetcher_contention.tex
%% Fig.4 — Prefetcher contention sequence diagram (단일 컬럼 fit, 단순화)

\begin{figure}[t]
\centering
\resizebox{0.95\columnwidth}{!}{%
\begin{tikzpicture}[
    x=1.0cm, y=0.55cm,
    actor/.style={font=\scriptsize\bfseries, align=center, fill=white, inner sep=2pt},
    lane/.style={dashed, gray!40},
    arrow/.style={-{Stealth[length=1.5mm]}, thick},
    msg/.style={font=\tiny, align=center, fill=white, inner sep=1pt},
    box/.style={rectangle, draw, minimum width=1.1cm, minimum height=0.40cm,
                font=\tiny, align=center, inner sep=1pt},
    label/.style={font=\scriptsize\itshape, align=center},
]
% Lanes
\node[actor] at (0, 4.7)  {CPU};
\node[actor] at (2, 4.7)  {Prefetcher};
\node[actor] at (4, 4.7)  {PCIe};
\node[actor] at (6, 4.7)  {GPU};
\foreach \x in {0, 2, 4, 6} {
    \draw[lane] (\x, 4.4) -- (\x, -1.3);
}

% Top scenario: Regular pattern
\node[label, text=red!60!black] at (3, 3.7) {Regular pattern $\rightarrow$ contention};
\node[box, fill=blue!10]   at (0, 3.0) {SIMD};
\draw[arrow] (0.3, 3.0) -- (1.7, 3.0);
\node[box, fill=yellow!25] at (2, 2.4) {active};
\draw[arrow] (2.3, 2.4) -- (3.7, 2.4);
\node[box, fill=red!22]    at (4, 1.8) {busy};
\draw[arrow] (4.3, 1.8) -- (5.7, 1.8);
\node[box, fill=red!12]    at (6, 1.2) {stall};

% Divider line
\draw[gray!30, thin] (-0.5, 0.6) -- (6.5, 0.6);

% Bottom scenario: Branchy pattern
\node[label, text=green!50!black] at (3, 0.1) {Branchy pattern $\rightarrow$ PCIe clear};
\node[box, fill=blue!10]   at (0, -0.5) {branchy};
\draw[arrow] (0.3, -0.5) -- (1.7, -0.5);
\node[box, fill=gray!25]   at (2, -0.5) {suppressed};
\node[box, fill=green!22]  at (4, -1.0) {free};
\draw[arrow] (4.3, -1.0) -- (5.7, -1.0);
\node[box, fill=green!12]  at (6, -1.0) {progress};
\end{tikzpicture}%
}
\caption{Regular 작업은 CPU prefetcher 의 활성을 유도하여 PCIe 버스에 speculative fetch 트래픽을 발생시켜 GPU 와 contention. Branchy 패턴 (rank/sort) 은 분기 예측 불가로 prefetcher 활성이 억제되어 PCIe 가 GPU 전용으로 유지된다.}
\label{fig:prefetcher-contention}
\end{figure}
